\centerline{\bf \Huge Муниципальные неравенства}


\problem{1}
$\bigl[\textbf{Питер 24/25, 10.2}\bigr]$ 
Известно, что $a>0, b>0, c>0$ и $a b+b c+a c>a+b+c$. Докажите, что $a+b+c>3$.

\problem{2}
$\bigl[\textbf{Новосибирск 23/24, 10.4}\bigr]$ 
Докажите, что для любых положительных чисел $a, b, c$, где $a \geq b \geq c$, выполняется неравенство:

$$
\sqrt{c(a-c)}+\sqrt{c(b-c)} \leq \sqrt{a b}
$$


\problem{3}
$\bigl[\textbf{Республика Татарстан, 22/23, 9.2}\bigr]$ 
Известно, что $2 x+y^2+z^2 \leqslant 2$. Докажите, что $x+y+z \leqslant 2$.

\problem{4}
$\bigl[\textbf{Саратовская область, 20/21, 9.3}\bigr]$ 
Докажите, что для любых положительных $p, q, r$ выполняется неравенство

$$
\left(\frac{1}{p q}+\frac{1}{q r}+\frac{1}{r p}\right)^2 \geq \frac{3}{p q r}\left(\frac{1}{p}+\frac{1}{q}+\frac{1}{r}\right) .
$$

\problem{5}
$\bigl[\textbf{Брянская область, 24/25, 11.1}\bigr]$ 
Докажите, что числа $x, y, z$ имеют один и тот же знак тогда и только тогда, когда одновременно выполняются следующие условия:

$$
\boldsymbol{x} \boldsymbol{y}+\boldsymbol{y} z+\boldsymbol{z} \boldsymbol{x}>\mathbf{0} \text { и } \quad \frac{1}{x y}+\frac{1}{y z}+\frac{1}{z x}>0 .
$$